\section{模型假设与符号说明}
\subsection{模型假设}
在题目给定的设备参数和交易规则基础上，作如下运行与信息假设，以明确电量流向、损耗计量及决策时可用的信息。

\textbf{（1）光伏优先直接供给小区负载。}光伏超过负载的部分可用于充电，外网购电既可直接供负载，也可给电池充电；三种电源允许共同供电。光伏优先是本文采用的运行约定。计划电量中能够即时使用或存储的部分予以接收，其余不接收，但普通购电仍按有效计划量结算。最终不能使用或储存的余电不计售电收入，也不对已支付的未利用电量重复计费。

\textbf{（2）采用固定充放电效率与外侧电量计量。}充电效率$\eta_{\mathrm c}$和放电效率$\eta_{\mathrm d}$均取0.90。充电量指进入储能设备前的电量，放电量指设备输出后可供负载使用的电量，5000 kW功率上限也按相应外侧计量。储电状态以电池内部电量表示，计算中不额外引入题目未给出的自放电和寿命折旧项。

\textbf{（3）以10分钟为调度区间。}问题一将附件1的每个采样值作为其后10分钟内的恒定功率和电价，以便将功率转换为时段电量。真实日期中的次日零点按实际日期衔接；问题一使用重复代表日，末尾零点标签映射到本周期开头。

\textbf{（4）储电状态连续衔接。}问题一主方案以6000 kWh开始，并在日末恢复到同一电量；不同起点的影响在本问灵敏度检验中单独考察。问题二至问题四承接前一天的剩余储电量，不逐日重置，也不将问题一的日首尾相等条件直接施加到后续每一天。

\textbf{（5）决策只使用当时可获得的信息。}问题二根据此前已经发生的负载、光伏及当前储电量制定计划，不使用附件3预报；问题三仅使用相应决策时刻已发布的预报。问题四未来实际价格不可提前获知，两条分支分别保留问题二和问题三的调整权限。问题二1月1日缺少更早历史，采用问题一基于附件1得到的购电结构启动，从6000 kWh连续运行，2月初电量由此前运行承接。

\textbf{（6）调整订单以零点原计划为统一结算基准。}对问题三及问题四的调整分支，最终有效普通订单减少部分退回原款，再收取该部分50\%的违约费用；增加部分按1.5倍价格计费，紧急购电另按5倍价格计费。多次修改时旧账冲回，重叠部分不重复收费，采用对应供电时段的价格。调整结果表记最终普通购电总量，零点计划和各次修改分别留存。

\subsection{符号说明}
以$d$表示日期，以$k$表示当天的10分钟区间。第$k$段为$[(k-1)\Delta t,k\Delta t)$，$k=1,\ldots,144$，$\Delta t=1/6$ h。$S_{d,0}$为当天零点储电量，$S_{d,k}$为第$k$段结束时的储电量，一天共对应145个状态边界。问题一只考虑一个代表日，以下模型中省略日期下标$d$。

功率、时段电量和库存分别表示供用电速率、该段累计电量及边界时刻的剩余电量。三者的主要符号与单位列于表\ref{tab:symbols}；决策量的电量流向在问题一模型中进一步说明。

{\small\renewcommand{\arraystretch}{1.18}
\begin{longtable}{>{\centering\arraybackslash}p{3.0cm}>{\raggedright\arraybackslash}p{8.2cm}>{\centering\arraybackslash}p{3.3cm}}
\caption{主要符号及其含义}\label{tab:symbols}\\
\toprule 符号 & 含义 & 单位或取值 \\ \midrule
\endfirsthead
\multicolumn{3}{c}{表\thetable\quad 主要符号及其含义（续）}\\
\toprule 符号 & 含义 & 单位或取值 \\ \midrule
\endhead
\bottomrule\endfoot
$d,\ k$ & 日期下标、日内时段索引 & $k=1,\ldots,144$\\
$\Delta t$ & 一个调度时段的长度 & $1/6$ h\\
$L_{d,k},\ P^{\mathrm{PV}}_{d,k}$ & 小区负载功率、光伏功率 & kW\\
$E^{\mathrm L}_{d,k},\ E^{\mathrm{PV}}_{d,k}$ & 负载电量、光伏电量；由相应功率乘$\Delta t$得到 & kWh\\
$p_{d,k}$ & 对应供电时段的电价 & 元/kWh\\
$m_{d,k}$ & 光伏直接供给负载的电量 & kWh\\
$n_{d,k},\ r_{d,k}$ & 光伏直供后的负载缺口、光伏富余电量 & kWh\\
$x_{d,k}$ & 外网购电中送入储能设备的电量 & kWh\\
$y_{d,k}$ & 送入储能设备的光伏电量 & kWh\\
$z_{d,k}$ & 储能设备输出给负载的电量 & kWh\\
$S_{d,k}$ & 第$k$段结束时的电池内部储电量 & kWh\\
$q_{d,k}$ & 外网购电总量；问题一中由电量分配关系确定 & kWh\\
$E^{\mathrm{ch}}_{d,k},\ E^{\mathrm{dis}}_{d,k}$ & 外侧总充电量$x_{d,k}+y_{d,k}$、总放电量$z_{d,k}$ & kWh\\
$w_{d,k}$ & 未被负载或电池利用的光伏电量 & kWh\\
$S_0$ & 问题一初始储电量，主方案取6000 & kWh\\
$S_{\mathrm{cap}}$ & 设备额定最大容量 & 12000 kWh\\
$S_{\min},\ S_{\max}$ & 允许运行的储电量下限、上限 & 1200、10800 kWh\\
$P_{\max}$ & 外侧充电或放电功率上限 & 5000 kW\\
$\eta_{\mathrm c},\ \eta_{\mathrm d}$ & 充电效率、放电效率 & 均为0.90\\
$C$ & 问题一全天普通购电费用 & 元\\
\end{longtable}}
